\documentclass[11pt,a4paper]{article}

% ===== Paquetes básicos =====
\usepackage[spanish]{babel}
\usepackage[utf8]{inputenc}
\usepackage[T1]{fontenc}
\usepackage{lmodern}
\usepackage{url}
\usepackage{hyperref}
\usepackage{geometry}
\usepackage{setspace}
\usepackage{listings}
\usepackage{xcolor}


\geometry{margin=2.5cm}
\setstretch{1.1}

% Configuración para código
%\lstset{
%	basicstyle=\ttfamily\small,
%	breaklines=true,
%	showstringspaces=false,
%	columns=fullflexible
%}


\lstdefinelanguage{JavaScript}{
	% Palabras clave básicas
	morekeywords=[1]{
		break, case, catch, class, const, continue, debugger, default, delete, do,
		else, export, extends, finally, for, function, if, import, in, instanceof,
		let, new, return, super, switch, this, throw, try, typeof, var, void, while,
		with, yield
	},
	% Literales y tipos
	morekeywords=[2]{
		false, null, true, undefined,
		boolean, number, string, symbol, bigint
	},
	% Objetos y funciones globales
	morekeywords=[3]{
		Array, Boolean, Date, Error, EvalError, Function, JSON, Map, Math,
		Number, Object, Promise, Proxy, RangeError, ReferenceError, RegExp, Set,
		String, SyntaxError, TypeError, URIError, WeakMap, WeakSet,
		console, window, document
	},
	sensitive=true,                % distingue mayúsculas
	% Comentarios
	morecomment=[l]{//},
	morecomment=[s]{/*}{*/},
	morecomment=[s]{/**}{*/},      % estilo JSDoc
	% Cadenas
	morestring=[b]',
	morestring=[b]",
	morestring=[b]`                % template strings
}[keywords, comments, strings]



\lstdefinestyle{jsstyle}{
	language=JavaScript,
	basicstyle=\ttfamily\small,
	keywordstyle=[1]\color{blue}\bfseries,
	keywordstyle=[2]\color{orange!80!black},
	keywordstyle=[3]\color{teal!70!black},
	stringstyle=\color{green!50!black},
	commentstyle=\color{gray},
	numberstyle=\tiny\color{gray},
	numbers=left,
	stepnumber=1,
	breaklines=true,
	showstringspaces=false,
	columns=fullflexible
}


\lstdefinestyle{terminalstyle}{
	basicstyle=\ttfamily\small,
	backgroundcolor=\color{black!5},
	keywordstyle=\color{black},
	commentstyle=\color{black},
	stringstyle=\color{black},
	showstringspaces=false,
	breaklines=true,
}



\title{Explorando un WSDL y comparándolo con una API REST}
\author{Juan Alvarado}
\date{}

\begin{document}
	
	\maketitle
	
	\section{Requisitos previos}
	\begin{itemize}
		\item Navegador web actualizado en Windows o macOS.
		\item 	Un editor de texto simple (Notepad, VS Code, TextEdit) instalado.
	\end{itemize}
	
	\section*{Intrucciones}
	
	\begin{itemize}
		\item Se proporciona la URL de un servicio SOAP público: URL del WSDL (para abrir en el navegador o descargar):
		\href{http://www.dneonline.com/calculator.asmx?wsdl}{http://www.dneonline.com/calculator.asmx?wsdl}
		\item 	Abrir el WSDL en el navegador o descargarlo y abrirlo en el editor de texto para ver su contenido XML.
		\item 	Localizar visualmente las secciones principales: tipos (definiciones de datos), mensajes, portType (operaciones abstractas), binding (detalles de protocolo) y service (endpoints concretos).
		\item 	Elegir una operación del WSDL y anotar: nombre de la operación, estructura de la petición, estructura de la respuesta y tipos de datos involucrados.
		\item 	También se proporciona la sigueinte liga con la documentación pública de una API REST sencilla: URL: \href{https://restful-api.dev}{https://restful-api.dev}.
		\item 	Identificar en la documentación REST los elementos equivalentes al contrato: rutas (URIs), métodos HTTP, parámetros de entrada, formato de respuesta, códigos de estado.
		\item En parejas, completar una tabla conceptual donde se mapeen los conceptos de WSDL a los de la API REST (por ejemplo, “operación” equivale a “endpoint + método HTTP”).
		\item 	Escribir un breve párrafo comparando la rigidez del contrato SOAP/WSDL frente a la flexibilidad percibida en la API REST; incluir al menos una ventaja y una desventaja de cada enfoque.
		\item Compartir algunas conclusiones en grupo para consolidar la idea de contratos estrictos frente a APIs modernas más ágiles.
	\end{itemize}
	
	
\end{document}
